% 输入均匀，S={(1,0),(2,1)}；两张真值表在已见处相同、未见处相反。
% 固定输出示例 h(3)=0,h(4)=1；风险和结论对任意二分类输出成立。
\begin{tikzpicture}[foundation picture]
 \path[use as bounding box] (0,10) rectangle (112,60);
 \node at (19,56) {$S=\{(1,0),(2,1)\}$};
 \foreach \x/\id/\label in {10/1/0,28/2/1}{
  \draw[FigureStroke,fill=FigureBlueFill,line width=1pt] (\x,47) circle (2.6);
  \node at (\x,47) {$\id$};
  \node[below] at (\x,42) {$y=\label$};
 }
 \node[foundation model,fill=white,minimum width=30mm] (algorithm) at (19,31) {$h=A(S)$};
 \draw[foundation flow] (19,39)--(algorithm.north);
 \node[foundation output,fill=white,minimum width=38mm,minimum height=10mm] (output) at (19,17)
  {$h(3)=0,\quad h(4)=1$};
 \draw[foundation flow] (algorithm.south)--(output.north);
 % 相同输出分支进入两个真值比较，隐藏真值不流入学习算法。
 \draw[FigureStroke,line width=1.1pt] (output.east)--(44,17)--(44,44);
 \draw[foundation flow] (44,44)--(54,44);
 \draw[foundation flow] (44,21)--(54,21);
 \node at (70,57) {$x=3$};
 \node at (97,57) {$x=4$};
 \foreach \row/\truth/\domain in {44/0/0,21/1/1}{
  \draw[foundation frame,fill=white] (54,\row-9) rectangle (112,\row+9);
  \node at (48,\row+7) {$\mathcal P_{\domain}$};
  \node at (70,\row+4) {$h(3)=0$};
  \node at (97,\row+4) {$h(4)=1$};
  \node at (70,\row-4) {$y=\truth$};
  \node at (97,\row-4) {$y=\truth$};
 }
 \node[text=FigureBlue!65!black] at (81,40) {$\checkmark$};
 \node[text=FigureRed!70!black] at (108,40) {$\times$};
 \node[text=FigureRed!70!black] at (81,17) {$\times$};
 \node[text=FigureBlue!65!black] at (108,17) {$\checkmark$};
\end{tikzpicture}%
